\chapter*{引言}
\addcontentsline{toc}{chapter}{引言}
\markboth{引言}{引言}

\section*{周期数学的工作定义}

\begin{definition}[周期数学（工作定义）]
周期数学是本书提出的一项\textbf{研究纲领}，它以\textbf{群作用生成的对称结构}为统一语言，以\textbf{平移、缩放、旋转}为基本生成变换的候选分类，以\textbf{闭合轨道}为核心几何图像，尝试对不同学科中的周期与分形结构进行统一数学描述。其研究对象是具有状态空间、演化规则和对称结构的系统。周期数学并不主张所有系统都具有周期性，而是致力于为\textbf{可能定义周期结构的系统}，提供一套统一的诊断与描述语言。\textbf{需要强调}：上述表述是工作定义与研究纲领，不是形式化数学定义；其严格版本需要在后续章节中逐条限定条件后建立。
\end{definition}

\section*{哲学启发与认知方法}

《道德经》的“道生一，一生二，二生三，三生万物”与《阴符经》的“观天之道，执天之行”为本书提供了哲学起点与认知方法的\textbf{启发}。前者描述了一个生成序列：从单一本原出发，经由分化的阶梯，产生多样性。在周期数学中，这一序列\textbf{可以类比于}数系扩张的阶梯 $\mathbb{N} \to \mathbb{Z} \to \mathbb{Q} \to \mathbb{R} \to \mathbb{C} \to \mathbb{H}$——每一步扩张都引入新的逆运算，每一步逆运算都引入新的对称性。这一类比是启发性的，不是数学定理。后者则提示了一条方法论路径：观察不变性，抽象守恒律，发现对称性，识别闭合轨道。\textbf{观天之道}可以理解为识别系统中的守恒量，\textbf{执天之行}可以理解为利用守恒律预测系统的未来状态。

\section*{莱布尼茨的普遍文字与《易经》的数学基础}

莱布尼茨在读到《易经》卦象后写道：“这个发现让我确信，伏羲不仅是中国哲学的创始人，而且是全人类最古老的智者之一。他所留下的符号系统，包含了一种普遍的文字——一种能够表达所有思想的字符。”\textbf{他的具体发现}——自然数二进制展开各位上的严格周期性，第 $0$、$1$、$k$ 位的周期分别为 $2$、$4$、$2^{k+1}$——以及“是否命名为 Periodus numerorum”的史料考证，已在\textbf{序言}中给出，此处不再重复。

在《易经》方向上的进一步延伸是：卦象体系的生成结构与二进制分数的圆周逼近相对应（见\textbf{序言}），且《易经》的“易与天地准”在周期数学中可以\textbf{部分地}获得数学意义：卦象体系与某些周期结构之间存在同构关系。“同声相应，同气相求”可以类比为频率的共振与相位的吸引；“无平不陂，无往不复”可以类比为闭合轨道的必然性；“方以类聚，物以群分”可以类比为群作用下的轨道分类。这些类比是\textbf{启发性的}，其数学内容需要在具体章节中逐一界定。

天干（$T=10$）与地支（$T=12$）的周期卷积生成六十甲子，八卦（$T=8$）与八卦（$T=8$）的嵌套生成六十四卦——周转圆嵌套规则统一了小周期生成大周期的过程：\textbf{外圈每前进一步，内圈旋转一整圈}。阴阳乘五行得十天干，四象乘三才得十二地支，八卦乘八卦得六十四卦。\textbf{这些对应关系在本书中被表述为代数同构或群作用等价，而非断言古代文献本身已以现代群表示论语言表述}。

\section*{为什么是平移、缩放与旋转}

周期数学将平移、缩放与旋转确立为三种基本生成变换，并非来自经验归纳，而是基于以下三个层面的\textbf{动机}。\textbf{需要说明}：以下三点是\textbf{选此三者的理由}，而非“任何连续变换必可分解为三者复合”的数学定理；后者是一个更强的、需要逐例验证的命题。

\begin{enumerate}[label=(\roman*)]
    \item \textbf{代数层面}：三种变换分别对应三种最基本的连续群——平移对应加法群 $(\mathbb{R}^n,+)$，缩放对应乘法群 $(\mathbb{R}^+,\times)$，旋转对应圆群 $U(1)$。\textbf{在适当的抽象层次上}，许多保持某种内禀结构不变的连续变换可以分解为这三种群的复合；但这一分解并非对所有连续群成立，需要逐例判定。
    \item \textbf{几何层面}：平移改变位置而不改变姿态，缩放改变尺度而不改变形状，旋转改变姿态而不改变位置。这三种操作\textbf{在欧氏几何与相似几何的框架内}恰好穷尽了不改变系统内禀结构的连续变换。
    \item \textbf{数系层面}：每一步数系扩张都引入新的逆运算——减法引入平移方向（对称中心 $0$），除法引入比例方向（对称中心 $1$），复根引入旋转方向（对称中心 $i$）。这一对应是\textbf{启发性的}，其严格数学内容属于域扩张与自同构群的研究范围。
\end{enumerate}

\section*{梯度、旋度、散度与三种对称性的对应}

在向量分析中，梯度 $\nabla$、旋度 $\nabla\times$、散度 $\nabla\cdot$ 是三个基本的微分算子。三者各自的结构如下（三者的严格定义、前提与边界条件见「矢量分析」章；本书不把它们与三类基本变换并列成对应关系）：

\begin{enumerate}[leftmargin=4em, label={(\arabic*)}]
    \item \textbf{梯度}：梯度把标量场映为向量场，指向函数增长最快的方向，其模长等于最大方向导数（需 $f \in C^1$ 且该点梯度非零）。
    \item \textbf{散度}：散度把向量场映为标量场，度量每单位体积的净流出率；在哈密顿系统中散度为零等价于相空间体积守恒（刘维尔定理）。
    \item \textbf{旋度}：旋度把向量场映为向量场，刻画每一点处的局部旋转趋势，其模长等于局部角速度的两倍（需 $\mathbf{F} \in C^1$）。
\end{enumerate}

\textbf{需要强调}：散度最标准的结构是局部源汇、体积变化率、连续性方程与 Liouville 型结果，而缩放在连续群层面的生成元是 $x\,d/dx$，属于另一个层次的对象；二者可以研究关系，但不能在没有严格桥梁之前并列成对应定律。本书因此不把“哪个算子对应哪种变换”作为命题陈述。

\section*{为什么要有守恒律}

守恒律在周期数学中不是附加的物理条件，而是\textbf{周期结构得以在保守系统中稳定存在的逻辑前提}。\textbf{需要强调}：以下论证适用于保守系统；对耗散系统，需引入广义守恒量或吸引子概念，将在后续章节中讨论。

\begin{enumerate}[label=(\roman*)]
    \item \textbf{变换的可逆性}：一个变换能够被反复施加而不破坏系统结构，\textbf{在保守系统中}当且仅当存在至少一个量在该变换下严格不变。没有守恒量的变换在保守框架下通常是不可逆的——轨道闭合在逻辑上难以保证。
    \item \textbf{轨道的约束性}：守恒律的数学本质是将系统演化约束在不变子流形上。没有这种约束，系统的演化是自由漂移，\textbf{一般}不产生可识别的规则轨道。\textbf{但在耗散系统中，吸引子可以替代守恒量起到约束作用}。
    \item \textbf{可观测的规律性}：守恒量是跨越时间的同一性——它让“过去的状态”与“现在的状态”之间建立起可判定的联系。
\end{enumerate}

从“周期基础”章建立的代数框架来看，周期系统完全可以只有平凡不变量（例如有限循环置换：状态集合有限、变换有限阶、每个状态都周期，但不一定存在非平凡守恒量）。因此\textbf{守恒量并非周期性的必要条件，也一般不是充分条件}。本书的基本原理表述为：

\textbf{在保守系统中，守恒量能够为轨道提供不变约束，从而为周期结构的存在与稳定性提供重要条件；但守恒量并非周期性的充分条件，也一般不是周期性的必要条件。在本书框架内，真正直接产生周期的代数枢纽是有限阶作用：}\[
\text{有限阶作用} \;\Longrightarrow\; \text{一致闭合轨道} \;\Longrightarrow\; \text{周期}.
\]

守恒与对称在其中的正当角色是：\textbf{发现、约束、解释周期结构的辅助工具}，而非逻辑链上不可绕过的环节。

\section*{为什么需要群表示论与流形框架}

\begin{enumerate}[label=(\roman*)]
    \item \textbf{群表示论回答“是什么变换”}：群表示论提供变换的代数结构——变换的集合是否构成群？群元素是否有有限阶？轨道是什么结构？
    \item \textbf{切空间回答“变换的瞬时方向”}：切空间捕捉系统在当前状态下的无穷小变化方向。一个闭合轨道在每一点处的切向量，构成了该周期运动的瞬时速度场。切空间中的向量场是否可积，决定了轨道能否在有限时间内闭合。
    \item \textbf{法空间回答“什么约束了变换”}：法空间描述系统在演化过程中受到的约束方向。守恒律在几何上表现为法空间中的约束超曲面——系统被限制在这个超曲面上运动，无法脱离。法空间的维度等于守恒量的个数，法空间的余维度等于系统的自由度。
    \item \textbf{角空间回答“变换进行到了哪里”}：角空间记录系统在闭合轨道上的相位位置。对于周期为 $T$ 的系统，角空间是圆群 $U(1)$ 或有限循环群 $\mathbb{Z}_T$ 的参数化。角空间的引入使得“相位”成为一个独立的几何维度，而不仅仅是时间的函数。\textbf{需要指出}：角空间的具体类型依赖于轨道的基本群与同伦类，并非所有闭合轨道都自动具有 $U(1)$ 或 $\mathbb{Z}_T$ 结构；其严格分类需要拓扑条件。
    \item \textbf{流形框架回答“全局结构是什么”}：流形是切空间、法空间与角空间的全局整合。一个周期系统的完整状态空间是一个流形，闭合轨道是这个流形上的闭曲线，切空间是流形在轨道各点处的切平面，法空间是切平面的正交补，角空间是轨道本身的参数化。
\end{enumerate}

\textbf{五者组合才是完整判定}：群表示论提供代数可能性，流形框架提供几何舞台。只有将群作用、切空间、法空间、角空间统一在流形框架中，才能同时回答：代数上群元素是否有有限阶？几何上轨道是否闭合？切方向上是否存在可积的向量场？法方向上是否存在守恒约束？角方向上相位是否单值？\textbf{这一框架是一个研究纲领}，其具体适用条件需要在各章中逐例建立。

\section*{流形框架的形式化}

设 $M$ 为系统的状态空间流形，$\gamma: [0,T] \to M$ 为一条闭合轨道，$p = \gamma(t)$ 为轨道上的一点。则：

\begin{enumerate}[leftmargin=4em, label={(\arabic*)}]
    \item \textbf{切空间} $T_p M$：轨道在 $p$ 点处的切向量 $\dot{\gamma}(t)$ 张成的一维子空间，描述系统的瞬时运动方向；
    \item \textbf{法空间} $N_p M$：切空间在 $T_p M$ 中的正交补，描述系统受到的约束方向。守恒量 $Q$ 的梯度 $\nabla Q(p)$ 是法空间中的向量；
    \item \textbf{角空间} $A_p$：轨道 $\gamma$ 的参数化空间，对于周期轨道，$A_p \cong U(1)$ 或 $A_p \cong \mathbb{Z}_T$，描述系统在轨道上的相位位置。\textbf{需要条件}：当轨道的基本群元素为无限阶且轨道可缩时，角空间自然为 $U(1)$；当基本群元素为有限阶时，角空间为 $\mathbb{Z}_T$。一般情形需要分别讨论。
    \item \textbf{切丛} $TM = \bigcup_{p \in M} T_p M$：所有切空间的集合，描述系统在所有可能状态下的运动方向；
    \item \textbf{法丛} $NM = \bigcup_{p \in M} N_p M$：所有法空间的集合，描述系统在所有可能状态下的约束方向；
    \item \textbf{角丛} $AM = \bigcup_{p \in M} A_p$：所有角空间的集合，描述系统在所有可能状态下的相位位置。
\end{enumerate}

\section*{守恒律与法空间的对应}

守恒律在流形框架中获得了几何解释：\textbf{在适当的非退化条件下}，一个守恒量 $Q$ 对应于法空间中的一个约束超曲面 $Q^{-1}(c)$。系统被限制在这个超曲面上运动，轨道 $\gamma$ 是约束超曲面上的闭曲线。切空间 $T_p M$ 中的向量与 $\nabla Q(p)$ 正交，这意味着守恒量在运动方向上不发生变化。法空间的维度等于独立守恒量的个数，法空间的余维度等于系统的有效自由度。\textbf{当独立守恒量的个数等于系统的自由度时}，轨道被完全约束为闭合曲线——这是完全可积系统的一个几何标志；\textbf{但反之不然}：闭合轨道的存在并不要求独立守恒量个数等于自由度。

\section*{角空间与相位算子的对应}

角空间是相位算子的几何实现。周期 $T$ 的第 $k$ 个相位对应复数 $\mathcal{O}_k^{(T)} = e^{i2\pi k/T}$，在角空间中表现为轨道上的均匀采样点；其三种作用（恒等、对径、旋转 $2\pi k/T$）与序言中相位算子的代数描述逐一对应，此处不再重复。角空间的拓扑结构决定了相位算子的代数结构：当角空间是连续圆 $U(1)$ 时，相位算子是连续旋转群；当角空间是有限循环群 $\mathbb{Z}_T$ 时，相位算子是有限旋转群。\textbf{需要指出}：上述对应以角空间的良定义性为前提，而角空间的严格构造依赖于轨道的同伦类与基本群，并非所有闭合轨道都自动满足。

\section*{周期卷积的初步定义}

传统的调和分析主要将信号分解为不同频率的正弦波之和，未必保留振幅与相位之间的交互信息。本书提出的一个核心研究纲领是周期卷积，其动机是将振幅（缩放）和相位（旋转）也纳入级数运算之中。

\begin{definition}[周期卷积：交互张量与合成映射（初步框架）]
设两个周期系统分别由周转圆级数
\[
f(x) = \sum_{n} A_n e^{i(\omega_n x + \phi_n)}, \quad g(x) = \sum_{m} B_m e^{i(\nu_m x + \psi_m)}
\]
表示（在适当的收敛条件下）。

\textbf{第一层（交互张量）}：定义两族周转圆分量之间的两两交互张量
\[
C_{n,m} = A_n B_m \operatorname{Similarity}^{(d)}(\theta_n, \theta_m),
\]
其中 $\theta_n = (\omega_n, \phi_n)$ 为 $f$ 的相位点，$\theta_m = (\nu_m, \psi_m)$ 为 $g$ 的相位点，$\operatorname{Similarity}^{(d)}$ 为 $d$ 维相位相似核。

\textbf{第二层（合成映射）}：定义从交互张量到函数空间的合成映射
\[
S(C)(x) = \sum_{n,m} C_{n,m} e^{i(\Omega_{n,m} x + \Phi_{n,m})},
\]
其中 $\Omega_{n,m}$ 为频率合成算子，$\Phi_{n,m}$ 为相位合成算子。

\textbf{周期卷积}（初步）定义为交互张量与合成映射的复合：
\[
(f \circledast g)(x) \;=\; \mathcal{K}^{(d)}\{f,g\} \;=\; S(C)(x).
\]
\end{definition}

\textbf{记号约定}：本书正文以 $\circledast$ 表示周期卷积算子；当需要显示相位相似核 $\operatorname{Similarity}^{(d)}$ 的阶数 $d$，或引入窗核 $W$ 时，同一个运算记作 $\mathcal{K}^{(d)}\{f,g\}$——封面即采用后者；应用各章中偶见简写 $A_{n,m} := A_n B_m$。

\textbf{需要强调}：上述定义是一个\textbf{研究纲领}，其严格成立需要明确以下条件：(i) 参与卷积的级数在适当的拓扑下收敛；(ii) 相位相似核 $\operatorname{Similarity}^{(d)}$ 满足对称性、有界性与归一化条件；(iii) 频率合成算子与相位合成算子有显式定义并满足结合律或分配律。这些条件将在后续章节中逐一给出或列为待研究问题。\textbf{在严格定义完成之前，上述框架应视为启发性纲领，而非已建立的数学对象。}

周期卷积以振幅（缩放）、频率（平移）、相位（旋转）、周期（定长）四个参数为算子，保留全部频率对与全部周期对的交互。这意味着，两个周期系统之间的相互作用，\textbf{在周期卷积的框架下}，不仅取决于它们各自的频率，还取决于它们的振幅之比、相位之差、以及周期之间的公倍关系。本书将尝试论证：经典积分变换——傅里叶变换、小波变换、梅林变换、拉普拉斯变换——\textbf{在适当的群作用、表示和函数空间框架下}可以在群表示与调和分析的框架下获得统一描述；\textbf{这比“它们本质上就是同一个东西”更为谨慎}——因为小波分析涉及仿射群、平移+缩放、表示、连续/离散小波、admissibility、具体测度等多个层次，精确的统一形式需要逐例检验。周期卷积的严格定义和性质将在后续章节中建立。

\section*{信息本体论与观察者框架}

\textbf{本节提出的信息本体论是本书的一个哲学立场，而非数学定理。}其数学内容仅在于：观察者的尺度与采样率决定了可被观察者数学化的信息范围。

设时间系统的合成向量为 $\mathbf{V}_T$，空间系统的合成向量为 $\mathbf{V}_S$，观察系统的合成向量为 $\mathbf{V}_P$。则
\begin{align}
\|\mathbf{V}_T\|^2 &= E \quad (\text{能量}), \label{eq:energy}\\
\|\mathbf{V}_S\|^2 &= m \quad (\text{质量}), \label{eq:mass}\\
\|\mathbf{V}_P\|^2 &= \sigma \quad (\text{尺度}). \label{eq:scale}
\end{align}
能量、质量、尺度——在此被\textbf{解释为}三个独立信息通量系统各自合成向量的模方。这一形式化的\textbf{数学内容}将在后续章节中给出；其\textbf{本体论解读}是哲学性的，不属于数学定理的范畴。

观察者自身也是信息结构的一部分。观察者的尺度 $\sigma$ 和采样率 $R(r,\sigma)$ 直接决定了可被数学化的信息范围。当采样点数 $N(r,\sigma) < 1$ 时，信息对观察者完全不可见；当采样率 $R$ 不足时，高频信息被混叠为低频伪影。\textbf{已被数学化的信息子集的边界，是观察者自身采样极限的函数}——这是一个关于观察者与信息关系的结构性陈述，其严格形式见「离散采样逼近、同余角与置闰」章 \cref{def:dynamic_sampling}、\cref{thm:invisibility_criterion}。

观察者采样窗口的选择是“信息数学化”的关键技术环节。当被观察的连续周期 $T_f$ 与观察者的采样间隔 $\Delta x$ 不可公度时，采样窗口长度的选择直接决定观察者可分辨的频谱质量。本书引入的\textbf{置闰方法}（见序言）为观察者提供了一个有严格最优判据的窗口选择机制：通过选取 $T_f/\Delta x$ 的连分数收敛子 $p_j/q_j$，观察者可在窗口长度 $T = p_j$ 上获得最小相位误差与最小频谱泄漏。这一机制使观察者框架从哲学陈述转化为可操作的算法准则。采样间隔与采样率的规范定义见「离散采样逼近、同余角与置闰」章 \cref{def:uniform_sampling_parallel,def:nyquist_sampling_parallel}。

\section*{信息的双重性}

\begin{enumerate}[leftmargin=4em, label={(\arabic*)}]
    \item \textbf{独立存在}：信息可以作为被辨识的单元而独立存在——一个群元素、一个函数值、一个周转圆半径、一个相位向量。
    \item \textbf{相互关联}：信息同样可以体现为实体之间的关系——群运算、函数复合、周转圆的嵌套叠加、相位向量的周期卷积投影。
\end{enumerate}

关联的实体化\textbf{可以}构成信息层级结构的基础：一阶信息（实体）、二阶信息（实体间的关系）、三阶信息（关系之间的关系）。\textbf{商群、函子、自然变换，可以被视为这一信息层次跃迁的若干数学实现}——这一对应是启发性的，其严格形式取决于具体的范畴论构造。

\section*{演化模型：守恒律的生住灭}

守恒量的时间导数
\[
\frac{dQ}{dt} = \{Q, \mathcal{H}\} + \frac{\partial Q}{\partial t}
\]
显式区分了内部动力学和外部驱动。生成（$\nabla Q \to 0$）、演化（$\nabla Q = 0$）、破缺（$\nabla Q \neq 0$）三阶段\textbf{在适当的辛几何框架下}可以获得严格表述。

\section*{研究目标与方法}

周期数学的研究目标，是建立一个能够描述不同动力系统\textbf{共同对称结构}的统一数学框架。我们关注的不是每个系统的具体物理、生物或社会机制，而是它们\textbf{共同的数学结构}。

周期数学接受现代数学的全部标准构造，不重新定义已有概念。其具体工作是：
\begin{enumerate}[leftmargin=4em, label={(\arabic*)}]
    \item 揭示对称性与方向性作为周期性的基础——三大基本对称（$0$、$1$、$i$）分别对应平移、缩放、旋转；
    \item 给出周期性在群表示论语言中的严格定义——周期函数是群作用下的不变函数，周转圆是旋转对称性的几何实现；
    \item 揭示数系扩张与对称群增加之间的结构对应；
    \item 统一纳入动态采样的观察者框架；
    \item 将积分变换\textbf{在适当条件下}统一纳入具有周期轨道的群上的调和分析框架；
    \item 以周期卷积统一描述跨周期相位运算——相位作为关系算子，周转圆嵌套生成大周期，六十干支是天干与地支的周期卷积合成；
    \item 建立\textbf{置闰方法}作为不可公度周期系统的通用逼近工具——以连分数最佳逼近为理论基础，以置闰误差有界定理为严格支撑，以置闰函子为范畴论实现，把周期分析中的频谱泄漏从不可控降为有界。
\end{enumerate}

\section*{与已有数学的关系}

周期数学不是要替代或修正群表示论、调和分析、表示论等已有学科，而是以“周期”为视角重新组织这些学科中与周期结构相关的部分。这一统一性的核心洞察是：
\begin{center}
\textbf{梯度、旋度、散度是向量分析中严格定义的微分算子，各自有其可微性前提与定义域条件；本书不对它们与三类基本变换作出对应断言，两组对象之间的任何关系都需逐例建立。}\\
\textbf{若选择适当的群作用、表示和函数空间，不同积分变换可以在群表示与调和分析的框架下获得统一描述。}\\
\textbf{能量、质量、尺度可以被解释为三个通量系统合成向量模方的观察者投影——这是一个研究纲领，而非已证定理。}\\
\textbf{置闰是用整数列逼近无理周期比的有界误差方法，其最优性由连分数收敛子给出——这是本书对周期分析方法论的一项独立贡献。}
\end{center}

本手稿定位为\textbf{研究纲领}：它提供了一种新的组织方式和提问角度，其最终成立取决于未来能否产生非平凡的新数学结果。

\section*{研究原则与步骤}

\begin{enumerate}[leftmargin=5em, label={\textbf{原则 \arabic*}}]
    \item \textbf{结构优先于机制}。关注系统之间的共同数学结构，而非系统特定的因果机制。
    \item \textbf{有限阶作用为核心}。在本书框架内，真正直接产生周期的代数枢纽是有限阶作用：有限阶 $\Rightarrow$ 一致闭合轨道 $\Rightarrow$ 周期。守恒与对称不是普遍的必要条件，而是发现与约束周期结构的辅助工具。
    \item \textbf{诊断而非断言}。提供工具来测量系统在各个对称变换下的闭合程度。
    \item \textbf{周期有生灭}。守恒律的维持保障周期的演化，耗散或外部驱动的介入导致周期的消亡。
\end{enumerate}

\begin{enumerate}[leftmargin=5em, label={\textbf{步骤 \arabic*}}]
    \item \textbf{数学内核构建}。严格定义从群作用、轨道闭合到守恒量的数学命题与定理。
    \item \textbf{无穷小探针}。以梯度、旋度、散度作为三种对称性的微分表达的候选工具。
    \item \textbf{范畴统一}。利用范畴论，将不同类型的周期结构形式化。
    \item \textbf{跨学科诊断}。使用这套工具，精确指出一个系统为何拥有或缺失某种周期结构。
\end{enumerate}

\section*{本书的主要结果}

\textbf{本节列出的是本书实际建立的结果}。每一条均按其被证明的强度陈述，并在方括号中给出所在章节；序言与以上各节的启发性图景与开放问题，此处\textbf{刻意不再重复}。

\begin{enumerate}[leftmargin=4em, label={(\arabic*)}]
    \item \textbf{两条线索}（第一章「两条线索：迭代生成运算与变换」）。在参数层面上，迭代把低阶运算组织为高阶运算，而反复施加的变换其参数沿同一层级上升。\textbf{运算线}为加法 $\to$ 乘法 $\to$ 幂（超运算层级）；\textbf{变换线}为平移 $\to$ 缩放 $\to$ 旋转；两条线索之间的关系是参数层面的对应，以定义与命题的形式给出，而非断言二者是同一回事。本书的三种生成变换由此获得生成性的推导，而不仅仅是一份枚举。
    \item \textbf{周期性的代数枢纽}（第一、十三至十五章）。在本书框架内，直接产生周期的代数枢纽是\textbf{有限阶作用}：有限阶 $\Rightarrow$ 一致闭合轨道 $\Rightarrow$ 周期。周期函数的周期集与具有周期轨道的群的结构，在「周期函数与具有周期轨道的群」章中展开；守恒与对称作为发现与约束周期结构的辅助工具进入，而非演绎链条上不可绕过的环节。
    \item \textbf{置闰误差有界与连分数最优性}（第二十一章）。对不可公度的周期比 $\rho$，下限、上限与均衡三种置闰序列均满足 $N_m \in \{\lfloor\rho\rfloor, \lceil\rho\rceil\}$，其累计误差始终有界——分别落在 $(-1,0]$、$[0,1)$ 与 $[-1/2,1/2]$ 之内；置闰窗口的最优选择由 $\rho$ 的连分数收敛子 $p_j/q_j$ 给出。本书把置闰方法作为对周期分析方法论的一项独立贡献：它把不可公度性的代价从不可控降为有界。
    \item \textbf{同余角与广义中国剩余定理}（第二十一章）。\textbf{同余角}把多个周期的相容性条件化为圆群上的角度条件，并把中国剩余定理从互素情形推广到可表述、可组合、可求解多周期耦合约束的形式；置闰由此从单周期的泄漏抑制升级为多周期的联合抑制。
    \item \textbf{周期卷积的两层结构}（第六、二十三、四十五章）。周期卷积定义为交互张量与合成映射的复合，$(f\circledast g)(x) = \mathcal{K}^{(d)}\{f,g\} = S(C)(x)$，使振幅（缩放）与相位（旋转）进入级数运算；其在平移周期范畴上的函子性以定理形式给出。
    \item \textbf{周转圆激活函数与函数逼近}（第四十二、四十五章）。在周转圆嵌套生成规则——外圈每前进一步，内圈旋转一整圈——的基础上，定义周转圆激活函数，并建立相应的函数逼近结果。本书的研究纲领在此取得可计算的具体形态。
\end{enumerate}

\textbf{本书不主张什么}：上述是本书已建立的结果。研究纲领中更强的命题——“任何保持某种内禀结构的连续变换都可分解为三种生成变换”“『哪个算子对应哪种变换』构成一条对应定律”“一般周期函数是有限个周转圆的叠加”——\textbf{均不在本书的主张之列}，正文中只以研究纲领或开放问题的形式承载。需要与之一并分清的是：在\textbf{复平面单参数保角变换族}这一有明确前设的框架内，三种基本变换的完备性是第一章\cref{sec:completeness_three_transformations}~节严格证明的结论；上述不主张的，是把它推广到「任意内禀结构的任意连续变换群」。仍未解决的问题集中陈述于第四十六章。

\section*{研究结构与跨学科应用}

本书分为三个部分，另以一章收束：
\begin{enumerate}[leftmargin=6em, label={\textbf{第 \arabic* 部分}}]
    \item \textbf{数学基础与生成元}（第一章至第十章）。从周期基础、数系、标量与矢量、函数与极限、实数分析、级数与卷积、复数分析、四元数分析、微分方程与解析到线性代数。
    \item \textbf{核心结构、工具与方法}（第十一章至第二十四章）。从流形基础与纤维丛、周期流形、周期范畴、概率论与分布、矢量分析，经周转圆组与周转球组、传播与演化、群与群函数、高阶群与高阶函数、周期函数与具有周期轨道的群，到离散采样逼近（含同余角）与置闰、调和分析、时序分析、多周期耦合系统。
    \item \textbf{跨学科应用与诊断}（第二十五章至第四十五章）。概率、密码、天文、历法、地理、易经、电磁、音频、机械、电路与芯片、生物、心理、社交、气候、经济、政治等系统的周期结构，以及数据分析、智能体、周转圆激活函数、算法系统、从数据到周转圆的建模。
\end{enumerate}
最后的\textbf{第四十六章「总结与展望」}汇聚全书三条主线，并陈述尚未解决的问题。

跨学科应用的理论必要性在于：一个统一框架的价值，在于它能否在多个独立领域中经受检验。周期数学不仅描述已知的周期系统，更诊断那些“看似有周期却未必满足闭合条件”的系统。行星轨道、晶体平移、基因振荡、科赫雪花、莱布尼茨二进制列周期、《易经》卦象体系——\textbf{在适当的抽象层次上}共享相同的群表示论结构和守恒律逻辑。

周期数学在每个应用章节中首先提问：\textbf{这个系统是否具有可识别的对称群？是否存在对应的守恒量？其状态空间中的演化轨道是否闭合？}（三个判据的哲学表述，以及“满足／部分满足／不满足”三种情形各自的处理方式，见\textbf{序言}「本书的立场与限度」。）

\section*{章节结构与阅读指南}

本书各章的体例与标注约定（“定义 $\to$ 命题 $\to$ … $\to$ 注释”的陈述顺序，以及“猜想／类比／研究纲领”三级标签的使用规则）见\textbf{序言}「体例与标注约定」。各章之间的关系可概括为三条主线：

\begin{enumerate}[leftmargin=4em, label={(\arabic*)}]
    \item \textbf{算子线}（梯度 $\nabla$ $\to$ 旋度 $\nabla\times$ $\to$ 散度 $\nabla\cdot$）。从「周期基础」章出发，经「矢量分析」章，贯穿「周期函数与具有周期轨道的群」章、「调和分析」章，逐一建立三个算子各自的可微性前提与适用范围。\textbf{这条线索不对三个算子与三类基本变换作出对应断言}；两组对象之间的任何关系均逐例讨论，「矢量分析」章明确声明两组之间不作对应断言。
    \item \textbf{数系线}（$\mathbb{N} \to \mathbb{Z} \to \mathbb{Q} \to \mathbb{R} \to \mathbb{C} \to \mathbb{H}$）。揭示周期结构如何随数系扩张而层层丰富。
    \item \textbf{观察者线}（尺度 $\sigma \to$ 采样率 $R \to$ 不可见性判据）。在“周期范畴”章中完成观察者与周期结构的范畴论统一。
\end{enumerate}

三条主线最终汇聚于一个核心命题：\textbf{在保守系统中，周期结构可能与守恒律密切相关，而守恒律在诺特定理框架下与对称性相关；但周期的真正代数枢纽是有限阶作用，守恒与对称是发现、约束、解释周期的辅助工具。}周期数学的任务，是以群作用为语言，以闭合轨道为判据，精确刻画从有限阶作用到周期性的演绎链条，同时明确守恒与对称在其中的恰当位置。

